\chapter{Refinement Background}

\section{Proactive Enforcement}

First we recall some definitions from Metric First Order Temporal Logic (MFOTL) and Lex \cite{whyenf, lex}.
Let $\Types$ be a set of base types and $\Variables$ a set of variables.

\begin{definition}\label{def:signature}
    A \emph{Signature}
    $\Sigma = \left( \left(\Domain_{ty}\right)_{ty \in \Types}, \Events, \Functions, \iota \right)$
    is a quadruple consisting of a domain of values such that
    for every $ty \in \Types$, $\Domain_{ty}$ is the set of values of type $ty$,
    a set of event names $\Events$, a set of function names $\Functions$,
    and an arity function $\iota : \Events \cup \Functions \to \Types^*$
    such that for every $e \in \Events$, $\iota(e) = (p_1,\dots,p_k)$ provides the
    parameter types of $e$
    and for any $f \in \Functions$, $\iota(f) = (a_1,\dots,a_k,r)$ provides the
    argument types $(a_1,\dots,a_k)$ and the return type $r$ of $f$.
\end{definition}

\begin{definition}\label{def:event}
    An \emph{Event} $e(p_1,\dots,p_k)$ consists of an event name
    $e \in \Events$ and parameters $(p_1,\dots,p_k)$ such that
    for every $1 \leq i \leq k$, $p_i \in \Domain_{\iota(e)_i}$.
\end{definition}

\begin{definition}\label{def:database}
    A \emph{Database} $D \in \Databases$ is a set of events.
\end{definition}

\begin{definition}\label{def:trace}
    A \emph{Trace} $\sigma \in \Traces$ is a sequence
    $\Sequence[1 \leq i \leq k]{(\tau_i, D_i)}$, $k \in \Nat \cup \lbrace \infty \rbrace$
    of timestamps $\tau_i \in \Nat$ and finite databases $D_i \in Databases$,
    where timestamps grow \emph{monotonically}
    ($\forall i < |\sigma|.\,\tau_i \leq \tau_{i+1}$)
    and \emph{progress}
    (if $|\sigma| = \infty$, then $\lim_{i\to\infty} \tau_i = \infty$).
\end{definition}

\begin{definition}\label{def:infinite-trace}
    The set of \emph{Infinite Traces} is denoted by $\Traces_\omega$.
\end{definition}

\begin{definition}\label{def:finite-trace}
    The set of \emph{Finite Traces} is denoted by $\Traces_f$.
\end{definition}

\begin{definition}\label{def:time-point}
    A \emph{time-point} $0 \leq i \leq |\sigma|$ is an index in a trace $\sigma$.
\end{definition}

\begin{definition}\label{def:empty-trace}
    The \emph{Empty Trace} is denoted by $\emptytrace$.
\end{definition}

\begin{definition}\label{def:property}
    A \emph{Property} is a subset $P\subseteq \Traces_\omega$.
\end{definition}

A Property can be induced by some policy $\phi$, we denote this policy by
$P(\phi)$.
We say a trace $\sigma$ is compliant with policy $\phi$ if and only if
$\sigma \in P(\phi)$.

\begin{definition}
   An \emph{Enforcer Function}
   $\mu_\phi : \Traces \times \Nat \times \Databases \to \Databases$
   for some policy $\phi$ takes as input a trace $\sigma \in \Traces$,
   a timestamp $\tau \in \Nat$, and a database $D \in \Databases$ and outputs a database $D' \in \Databases$
   such that if the trace $\sigma \in P(\phi)$,
   $\sigma \cdot \left(\tau, D'\right) \in P(\phi)$.
\end{definition}
% TODO: elaborate on functionality of a proactive enforcer

\section{Enforcement as an Event System}

We first define enforcement as an event system and then show how the notion of refinement as in Sprenger et. al \cite{igloo} works for such an event system.

\begin{definition}
    \label{def:enf-state-space}
    The \emph{Enforcement State Space} $\Statespace = \Nat \times \Traces$
    contains tuples $(\tau,\sigma)$.
\end{definition}
The timestamp $\tau$ can be seen as the \emph{current} timestamp
and $\sigma$ the trace up until timestamp $\tau$.

\begin{definition}
    \label{def:enf-initial-states}
    The \emph{Initial States}
    $\Initialstates = \lbrace (\tau, \emptytrace) | \tau \in \Nat \rbrace$
    are all states with an arbitrary timestamp $\tau$ and an empty trace.
\end{definition}
Intuitively this can be interpreted as that an enforced system may be
started at an arbitrary timestamp.

\begin{definition}
    \label{def:enf-event-space}
    The \emph{Event Space} $\Eventspace=\Databases$ is 
\end{definition}

\begin{definition}
    \label{def:enf-transition-function}
    The \emph{Transition Function} $\Transitionfunction{\phi}{\text{hello world}} : $
\end{definition}